\documentclass[10pt,a4paper]{amsart}
\usepackage[margin=19mm]{geometry}
\usepackage{amsmath,amssymb,mathtools}
\usepackage[scr=rsfs,cal=euler]{mathalfa}
\usepackage{xfrac}
\usepackage[unicode,hidelinks,bookmarks=false]{hyperref}
\newcommand{\ie}{i.e.\ }
\newcommand{\cf}{cf.\ }
\newcommand{\resp}{respectively\ }
\newcommand{\Subsection}{\subsection}
\providecommand{\nobreakdash}{\nobreak\hbox{-}}
\setlength{\emergencystretch}{2em}
\multlinegap0pt
\usepackage{amssymb}
\usepackage{dsfont}
\usepackage{enumerate}
\usepackage{tikz}
\usepackage{mathtools}
\usepackage{cancel}
\usepackage[shortlabels]{enumitem}
\newtheorem{theorem}{Theorem}
\usepackage{cleveref}
\crefname{theorem}{Theorem}{Theorems}
\newcommand{\Gc}{\mathcal{G}}
\newcommand{\PNc}{\mathcal{PN}}
\DeclareMathOperator*{\argmax}{arg\:max}
\newcommand{\boundary}{B}
\newcommand{\edgelength}{\omega}
\newcommand{\edgeset}{E}
\newcommand{\inflap}{\Delta_{\infty}}
\newcommand{\internalnodes}{\nodeset\setminus\boundary}
\newcommand{\nodeset}{V}
\title{Nonlinear spectral graph theory}
\author{Piero Deidda, Francesco Tudisco, Dong Zhang}
\date{Source: arXiv:2504.03566v2 (CC BY 4.0)}
\begin{document}
\maketitle

\noindent\emph{Source and licence.} Nonlinear spectral graph theory. Version arXiv:2504.03566v2 (CC BY 4.0), distributed by arXiv under the Creative Commons Attribution 4.0 International licence, http://creativecommons.org/licenses/by/4.0/, as recorded in the arXiv licence metadata for this exact version and shown on the abstract page https://arxiv.org/abs/2504.03566v2. Full licence notice: \texttt{licenses/arxiv-2504.03566-CC-BY-4.0.txt}.\par\smallskip

\noindent\textit{Editorial scope.} Counted unit: the article's theorem Thm\_spectral\_min\_partitions (source0.tex lines 4282--4294). The article's complete original proof of it is delivered with it (source0.tex lines 4296--4349, one \texttt{proof} environment of 54 lines). No mathematics is written, repaired or re-derived: the statement and the proof are the article's own lines, in the article's own notation, and no retained line is substituted or re-flowed.  Retained uncounted as the source-local context the proof needs: lines 3586--3590.  Nothing was omitted from any retained span, so the package carries no omission record.  No retained line contains a citation, so the wrapper carries no bibliography and no bibliography entry.  No retained line contains a figure environment, an image-inclusion command or drawing code, so no image is delivered and the package declares no asset dependency.  Editorial formatting only, in addition: an AMS \texttt{amsart} preamble that restates the article's own declarations and macros used by the retained lines, namely the environments \texttt{theorem} on separate counters, together with the standard packages those lines need.  The retained environments print in this document as Theorem 1 in order of appearance, whereas the article prints the same items with the numbering of its own full document.  The complete native source of the article as distributed in the arXiv e-print for this exact version is bundled unchanged as \texttt{sources/175989/source0.tex}.
\begin{theorem}\label{thm:k-inequality}
Given a connected graph $\Gc$, we have the equalities  $\Lambda_1(\inflap)=1/R_1(\Gc)$ and $\Lambda_2(\inflap)=1/R_2(\Gc)$. Moreover, for any $k=1,2,\cdots,N$,
$$\frac{1}{R_{\PNc(f_k)}(\Gc)}\le \Lambda_k(\inflap)\le \frac{1}{R_k(\Gc)}$$ 
 where $f_k$ is any eigenfunction corresponding to $\Lambda_k(\inflap)$ and $\PNc(f_k)$ is the number of perfect nodal domains induced by $f_k$. 
\end{theorem}

\begin{theorem}\label{Thm_spectral_min_partitions}
Let $\edgelength_M:=\min_{(u,v)\in\edgeset} \edgelength_{(u,v)}$ be the reciprocal of the length of the longest edge of a graph $\Gc$. Then the $\infty$-Laplacian spectral minimal partition costs and the packing radii of $\Gc$ satisfy the following inequalities:  
%
    $$
    R_k(\Gc) 
 \leq \frac{1}{\widehat{\Lambda}_k(\inflap)}\leq R_k(\Gc)+\frac{1}{2\edgelength_M} 
    %
    \qquad \text{and} \qquad 
    %
 R_k(\Gc)-\frac{1}{2\edgelength_M}\leq\frac{1}{\widetilde{\Lambda}_k(\inflap)}\leq R_k(\Gc).
    $$
    %
\end{theorem}

\begin{proof}
    Start by recalling the definition of $R_k(\Gc)$:
    %
    \begin{equation}
    R_k(\Gc)=\max_{v_1,\dots,v_k\in \internalnodes}\min_{i,j=1,\dots,k} \frac{d(v_i,v_j)}{2}\,.
    \end{equation}
    %
Then let $v_1,\dots,v_k$ be a maximizing set of nodes in the definition of $R_k(\Gc)$, and for any $i=1,\dots,k$, let
%
\begin{equation}
\widehat{V}_i=\{u\in \nodeset\,|\;d(u,v_i)< R_k(\Gc)\} \quad \text{and} \quad \widetilde{V}_i=\{u\in \nodeset\,|\;d(u,v_i)< R_k(\Gc)-1/(2\edgelength_M)\}\,.
\end{equation}
%
By construction, the families of sets $\{\widehat{V}_i\}$ and $\{\widetilde{V}_i\}$ are suitable sets in the definitions, respectively, of $\widehat{\Lambda}_k(\inflap)$ and $\widetilde{\Lambda}_k(\inflap)$. Hence, from the characterization of the first $\infty$-eigenvalue, see \Cref{thm:k-inequality}, we observe
%
\begin{equation}
\frac{1}{\widehat{\Lambda}_k(\inflap)}\geq \min_i \frac{1}{\Lambda_1(\inflap,\widehat{V}_i)}=\min_i \max_{u\in \widehat{V}_i}d_{\nodeset\setminus \widehat{V}_i}(u)\geq R_k(\Gc)\,.
\end{equation}
%
%
Analogously we observe that, for any $i$, 
%
\begin{equation}
\frac{1}{\widetilde{\Lambda}_k(\inflap)}\geq  R_k(\Gc)- \frac{1}{2\edgelength_M}\,.
\end{equation}
%

We miss to prove the upper bounds, let $V_1,\dots V_k$ be a maximizing family of sets in the definition of $\widehat{\Lambda}_k(\inflap)$\,, i.e.
%
\begin{equation}
\widehat{\Lambda}_k(\inflap)=\max_i \frac{1}{\max_{u\in V_i}d_{\nodeset\setminus V_i}(u)}\,,
\end{equation}
%
then define $v_i^{max}=\argmax_{u\in V_i}d_{\nodeset\setminus V_i}(u)$ and $\partial V_i:=\{u\in V_i\;| \;\exists (u,v)\in E\; \text{with}\; v\in \nodeset\setminus V_i\}$\,.
Note that any path connecting $v_i^{max}$ to $v_j^{max}$ has to be incident both to $\partial V_i$ and $\partial V_j$, thus
%
\begin{align}\label{eq:1_Thm_spectral_min_partitions_a}
d(v_i^{max},v_j^{max})\geq d(v_i^{max},\partial V_j) +d(\partial V_j,v_j^{max})\geq \frac{1}{\widehat{\Lambda}_k(\inflap)}+\frac{1}{\widehat{\Lambda}_k(\inflap)}-\frac{1}{\edgelength_M}
\end{align}
%
where we have used that for any $i$,
$d(v_i^{max},\nodeset\setminus V_i)\leq d(v_i^{max},\partial V_i)+1/\edgelength_M$ and $d(v_i^{max},\partial V_j)\geq d_{\nodeset\setminus V_i}(v_i^{max}) $.
Clearly \eqref{eq:1_Thm_spectral_min_partitions_a} implies that 
%
\begin{equation}
 R_k(\Gc)\geq \frac{1}{\widehat{\Lambda}_k(\inflap)}-\frac{1}{2\edgelength_M}.  
\end{equation}
%
The same proof, assuming $V_1,\dots, V_k$ to be a maximizing family of subsets in the definition of $\widetilde{\Lambda}_k$, leads to  
\begin{equation}
    R_k(\Gc)\geq \frac{1}{\widetilde{\Lambda}_k(\inflap)},
\end{equation}
in this case, in fact, the subsets $\{V_i\}$ are non-adjacent and thus $d(v_i^{max},v_j^{max})\geq d(v_i^{max},\nodeset\setminus V_i)+ d(v_j^{max},\nodeset\setminus V_j)$.
\end{proof}

\end{document}
